%%%%%%%%%%%%%%%%%%%%%%%%%%%%%%%%%%%%%%%%%%%%%%
%% preamble.tex --- 版式层 (LAYOUT LAYER)
%%
%% 这是唯一决定"长什么样"的文件：class、边距、段落、配色、
%% 定理环境的外观、表格的线。改这一个文件，所有章节一起变；
%% 各章正文 .tex 里一个版式指令都没有。
%%
%% 骨架来自 Harry 的作业模板；作业专用的 \problem / \answer
%% 计数器换成了笔记要用的编号定理环境。
%%
%% 正文层依赖的接口（替换本文件时必须继续提供）：
%%   环境 definition / claim / example  编号，共用一个计数器
%%   环境 note                          不编号
%%   环境 proof                         amsthm，以 ∎ 收尾
%%   命令 \fig{stem}{caption}{label}    按主题选 -light / -dark 图
%%   命令 \chapref{显示文字}{stem}       指向同目录另一章的 PDF
%%   命令 \secref{label}                § 号交叉引用
%%   环境 checklist + \checkitem        末尾的复选框清单
%%   开关 \ifdark                       由 build.sh 通过 \DARKMODE 置位
%%%%%%%%%%%%%%%%%%%%%%%%%%%%%%%%%%%%%%%%%%%%%%

\documentclass[12pt]{article}


%%%%%%%%%%%%%%%%%%%%%%%%%%%%%%%%%%%%%%%%%%%%
%% Page formatting.
%%%%%%%%%%%%%%%%%%%%%%%%%%%%%%%%%%%%%%%%%%%%%%
\usepackage[margin=1in]{geometry} % adjust margins
\setlength\parindent{0em} % never indent a paragraph
\setlength{\parskip}{1em} %Seperate paragraphs


%%%%%%%%%%%%%%%%%%%%%%%%%%%%%%%%%%%%%%%%%%%%
%% Packages
%%%%%%%%%%%%%%%%%%%%%%%%%%%%%%%%%%%%%%%%%%%%%%
\usepackage{amsthm, amssymb, amsmath} % basic math stuff -- always include this
\usepackage{enumerate}   % customize enumeration
\usepackage[table]{xcolor} % named colours (+ \arrayrulecolor), before hyperref
\usepackage{tcolorbox}   % allows me to make colored boxes
\usepackage{graphicx}    % the chapter figures
\usepackage{tabularx}    % tables whose text column fills the line
\usepackage{xltabular}   % ...and that may break across pages
\usepackage{array}       % >{...} column prefixes
\usepackage{caption}


%%%%%%%%%%%%%%%%%%%%%%%%%%%%%%%%%%%%%%%%%%%%%%
%% Theme. build.sh defines \DARKMODE for the dark build and
%% nothing for the light one. Page colour, link colour and the
%% table rules all have to agree, or a dark PNG lands in a white
%% margin (and vice versa).
%%%%%%%%%%%%%%%%%%%%%%%%%%%%%%%%%%%%%%%%%%%%%%
\newif\ifdark
\ifdefined\DARKMODE \darktrue \fi

\ifdark
  \definecolor{pagebg}{HTML}{14161A}
  \definecolor{ink}{HTML}{E6E8EB}
  \definecolor{linkcol}{HTML}{6FB7E8}
  \definecolor{rulecol}{HTML}{4A5057}
  \definecolor{boxbg}{HTML}{1E2126}
\else
  \definecolor{pagebg}{HTML}{FFFFFF}
  \definecolor{ink}{HTML}{000000}
  \definecolor{linkcol}{HTML}{0000FF}
  \definecolor{rulecol}{HTML}{000000}
  \definecolor{boxbg}{HTML}{F2F4F7}
\fi

\usepackage[colorlinks=true, allcolors=linkcol]{hyperref} % clickable hyperlinks

\ifdark
  \usepackage{pagecolor}
  \pagecolor{pagebg}
  \color{ink}
  \arrayrulecolor{rulecol}
\fi


%%%%%%%%%%%%%%%%%%%%%%%%%%%%%%%%%%%%%%%%%%%%%%
%% Numbered theorem environments.
%%
%% The head is run-in and bold, so a Definition still opens the
%% paragraph the way "\textbf{Problem 1.}" does in the homework
%% template --- it just carries a number now, and can be \ref'd.
%% definition / claim / example share one counter, so the numbers
%% read in document order and never collide with the § numbers.
%%%%%%%%%%%%%%%%%%%%%%%%%%%%%%%%%%%%%%%%%%%%%%
\newtheoremstyle{notes}
  {\topsep}    % space above
  {\topsep}    % space below
  {}           % body font
  {0pt}        % indent
  {\bfseries}  % head font
  {.}          % punctuation after head
  { }          % space after head -- a plain space makes it run-in
  {}           % head spec -- empty means the default layout
\theoremstyle{notes}
\newtheorem{definition}{Definition}
\newtheorem{claim}[definition]{Claim}
\newtheorem{example}[definition]{Example}

% Note is a caveat, not a result: unnumbered, nothing refers back to it.
\newtheorem*{note}{Note}

% amsthm sets the proof head in italic; this template wants it bold.
\renewcommand{\proofname}{\normalfont\bfseries Proof}


%%%%%%%%%%%%%%%%%%%%%%%%%%%%%%%%%%%%%%%%%%%%
%% Define (or redefine) any commands/macros here
%%%%%%%%%%%%%%%%%%%%%%%%%%%%%%%%%%%%%%%%%%%%%%
\newcommand{\R}{\mathbb{R}} % sets up the command \R to write the R symbol for real numbers
\newcommand{\Z}{\mathbb{Z}} % use \Z for integers
\newcommand{\Q}{\mathbb{Q}} % use \Q for rationals
\newcommand{\N}{\mathbb{N}} % use \N for natural numbers

% --- notes-specific -------------------------------------------------

% \fig{stem}{caption}{label} -- picks figures/<stem>-light.png or -dark.png
% to match the build, so the figure and the page always agree.
\graphicspath{{../figures/}}
\newcommand{\figvariant}{light}
\ifdark \renewcommand{\figvariant}{dark} \fi
\newcommand{\fig}[3]{%
  \begin{figure}[htbp]
    \centering
    \includegraphics[width=\linewidth]{#1-\figvariant}
    \caption{#2}\label{#3}
  \end{figure}}

% \chapref{text}{stem} -- link to a sibling chapter's PDF. Each chapter
% compiles to its own file, so cross-chapter links point at .pdf, and
% degrade to plain readable text if the sibling is not built yet.
\newcommand{\chapref}[2]{\href{#2.pdf}{#1}}

% \secref{label} -- "§ 4.3", clickable. This replaces the hand-typed
% section numbers the Markdown had to carry.
\newcommand{\secref}[1]{\S\ref{#1}}

% A forward pointer to another chapter or to code.
\newcommand{\pointer}[1]{\par\noindent$\rightarrow$~#1\par}

% The closing checklist.
\newenvironment{checklist}{\begin{itemize}\setlength{\itemsep}{0pt}}{\end{itemize}}
\newcommand{\checkitem}{\item[$\square$]}

% Interpunct, for titles like "05 · Modelling".
\newcommand{\midot}{\textperiodcentered}

%%%%%%%%%%%%%%%%%%%%%%%%%%%%%%%%%%%%%%%%%%%%%%

%%%%%%%%%%%%%%%%%%%%%%%%%%%%%%%%%%%%%%%%%%%%%%
%% Two consequences of \parskip{1em}, fixed.
%%%%%%%%%%%%%%%%%%%%%%%%%%%%%%%%%%%%%%%%%%%%%%

% The table of contents is a list, not a run of paragraphs; at 1em of
% skip per entry it spills onto a second page. Zero the skip for its
% duration only, so the body keeps the spacing the template wants.
\makeatletter
\let\notes@toc\tableofcontents
\renewcommand{\tableofcontents}{%
  \begingroup\setlength{\parskip}{0pt}\notes@toc\endgroup}
\makeatother

% \parskip does not reach inside a tabular, so rows come out tighter
% than the surrounding text. Open them back up.
\renewcommand{\arraystretch}{1.25}
%%%%%%%%%%%%%%%%%%%%%%%%%%%%%%%%%%%%%%%%%%%%%%
